\documentclass[12pt,letterpaper]{article}

\usepackage[utf8]{inputenc}
\usepackage[T1]{fontenc}
\usepackage[spanish,es-nodecimaldot,es-noshorthands]{babel}
\usepackage{geometry}
\usepackage{amsmath}
\usepackage{tikz}
\usepackage{xcolor}
\usepackage{float}
\usepackage{enumitem}
\usepackage{titlesec}
\usepackage{graphicx}
\usepackage{tabularx}
\usepackage{booktabs}
\usepackage{array}

\usetikzlibrary{
    arrows.meta,
    positioning,
    shapes.geometric,
    calc
}

\geometry{
    left=2.3cm,
    right=2.3cm,
    top=2.2cm,
    bottom=2.2cm
}

\titleformat{\section}{\large\bfseries}{}{0em}{}

\setlength{\parindent}{0pt}
\setlength{\parskip}{0.5em}

\begin{document}

% ---------------------------------------------------------
% TITULO
% ---------------------------------------------------------

\begin{center}
    {\LARGE \textbf{Proyecto Final: Control de un UR3 mediante gemelo digital y realidad mixta}}\\[0.3cm]
    {\large Selección de hardware, protocolos y toma de decisiones}
\end{center}

% ---------------------------------------------------------
\section{Propuesta}
% ---------------------------------------------------------

El objetivo del proyecto es desarrollar un sistema en el que un 
\textbf{robot colaborativo UR3} pueda ser controlado a partir de la 
manipulación de su \textbf{gemelo digital en realidad mixta}.

El operador utilizará unos lentes de realidad mixta para observar una
representación digital del robot y de un almacén previamente parametrizado.

El almacén será representado mediante una cuadrícula tridimensional. 
Por ejemplo, una estructura de:

\[
10 \times 10 \times 10
\]

permitiría representar hasta:

\[
10 \cdot 10 \cdot 10 = 1000
\]

posiciones lógicas diferentes.

Cada posición puede identificarse mediante una coordenada:

\[
P(i,j,k)
\]

donde $i$, $j$ y $k$ representan la posición del producto dentro del
almacén virtual.

El usuario podrá seleccionar dentro de los lentes una celda o producto.
El sistema calculará primero una trayectoria para el gemelo digital del
UR3 y comprobará que el movimiento sea válido.

Si la trayectoria no presenta colisiones, límites articulares o
singularidades, el usuario podrá presionar un botón de
\textbf{Ejecutar} dentro de la interfaz de realidad mixta.

Solamente después de una segunda validación, la trayectoria será enviada
al robot UR3 físico.

% ---------------------------------------------------------
\section{Arquitectura seleccionada}
% ---------------------------------------------------------

La arquitectura se divide en cuatro partes principales:

\begin{enumerate}[leftmargin=1.2cm]

    \item \textbf{Interfaz de realidad mixta:}
    Meta Quest 3 ejecutando una aplicación desarrollada en Unity y OpenXR.

    \item \textbf{Computadora de borde:}
    encargada de ROS 2, MoveIt 2, planeación de trayectorias,
    inteligencia artificial y procesamiento de imágenes.

    \item \textbf{Sistema físico:}
    robot UR3, gripper y cámara RGB-D colocada cerca del efector final.

    \item \textbf{Sistema de inventario:}
    sensores en el almacén, ESP32, Raspberry Pi, base de datos local,
    nube y aplicación móvil.

\end{enumerate}

% ---------------------------------------------------------
\section{Hardware seleccionado}
% ---------------------------------------------------------

\begin{table}[H]
\centering
\small
\renewcommand{\arraystretch}{1.35}

\begin{tabularx}{\textwidth}{
    >{\raggedright\arraybackslash}p{3.3cm}
    >{\raggedright\arraybackslash}p{3.6cm}
    X
}

\toprule
\textbf{Elemento} & \textbf{Hardware elegido} & \textbf{Función} \\
\midrule

Robot &
UR3 &
Ejecutar físicamente las trayectorias previamente validadas en el gemelo digital. \\

Realidad mixta &
Meta Quest 3 &
Visualizar el gemelo digital, seleccionar productos o posiciones y autorizar movimientos. \\

Procesamiento principal &
PC de borde con GPU NVIDIA &
Ejecutar ROS 2, MoveIt 2, YOLO, planeación de movimiento y agente de IA. \\

Visión del robot &
Cámara RGB-D RealSense D405 &
Identificar productos y obtener información de profundidad cerca del efector final. \\

Manipulación &
Gripper eléctrico ligero &
Tomar y depositar los productos del almacén. \\

Inventario &
Celdas de carga + HX711 &
Determinar la cantidad aproximada de productos mediante el peso almacenado en cada celda. \\

Adquisición de sensores &
ESP32 &
Leer los sensores de las diferentes secciones del almacén. \\

Servidor local &
Raspberry Pi 5 + SSD &
Ejecutar MQTT, API local y base de datos de inventario. \\

Red &
Switch Ethernet + punto de acceso Wi-Fi &
Comunicar robot, PC, Raspberry Pi y lentes dentro de la red local. \\

Nube &
Base de datos PostgreSQL en la nube &
Mantener una copia remota del estado del almacén y proporcionar información a la aplicación móvil. \\

\bottomrule
\end{tabularx}

\caption{Hardware principal seleccionado para el proyecto.}

\end{table}

% ---------------------------------------------------------
\section{Protocolos de comunicación}
% ---------------------------------------------------------

Cada tramo del sistema utiliza un protocolo diferente dependiendo de
las necesidades de velocidad, confiabilidad y cantidad de información.

\begin{table}[H]
\centering
\small
\renewcommand{\arraystretch}{1.4}

\begin{tabularx}{\textwidth}{
    >{\raggedright\arraybackslash}p{3.1cm}
    >{\raggedright\arraybackslash}p{3.3cm}
    >{\raggedright\arraybackslash}p{2.8cm}
    X
}

\toprule
\textbf{Origen} &
\textbf{Destino} &
\textbf{Protocolo} &
\textbf{Justificación} \\
\midrule

Meta Quest 3 &
PC de borde &
WebSocket / TCP-IP &
Permite enviar coordenadas, comandos y recibir en tiempo real el estado del gemelo digital. \\

PC de borde &
UR3 &
ROS 2 + RTDE sobre Ethernet &
Permite utilizar el driver oficial de Universal Robots y enviar trayectorias previamente calculadas. \\

Cámara RGB-D &
PC de borde &
USB 3.0 &
Permite transmitir imágenes RGB y mapas de profundidad con baja latencia. \\

ESP32 &
Raspberry Pi &
MQTT &
El inventario genera pequeños mensajes de telemetría y no necesita comunicación de tiempo real estricto. \\

Raspberry Pi &
PC de borde &
MQTT / REST &
Permite compartir información del inventario con el sistema robótico. \\

Raspberry Pi &
Nube &
HTTPS / TLS &
Sincroniza de manera segura los movimientos e inventario del almacén. \\

Aplicación móvil &
Nube &
HTTPS + WebSocket &
Permite consultar inventario y recibir actualizaciones sin conectarse directamente al robot. \\

\bottomrule
\end{tabularx}

\caption{Protocolos seleccionados para cada tramo del sistema.}

\end{table}

% ---------------------------------------------------------
\section{Diagrama general del sistema}
% ---------------------------------------------------------

\begin{figure}[H]
\centering

\resizebox{\textwidth}{!}{
\begin{tikzpicture}[
    node distance=1.2cm and 1.5cm,

    vr/.style={
        rectangle,
        rounded corners,
        draw,
        thick,
        minimum width=3.8cm,
        minimum height=1.2cm,
        align=center,
        fill=purple!12
    },

    control/.style={
        rectangle,
        rounded corners,
        draw,
        thick,
        minimum width=4.4cm,
        minimum height=1.3cm,
        align=center,
        fill=orange!15
    },

    robot/.style={
        rectangle,
        rounded corners,
        draw,
        thick,
        minimum width=3.8cm,
        minimum height=1.2cm,
        align=center,
        fill=blue!12
    },

    sensor/.style={
        rectangle,
        rounded corners,
        draw,
        thick,
        minimum width=3.8cm,
        minimum height=1.2cm,
        align=center,
        fill=green!12
    },

    data/.style={
        rectangle,
        rounded corners,
        draw,
        thick,
        minimum width=4cm,
        minimum height=1.2cm,
        align=center,
        fill=cyan!12
    },

    flecha/.style={
        -{Latex[length=2.5mm]},
        thick
    },

    doble/.style={
        {Latex[length=2.5mm]}-{Latex[length=2.5mm]},
        thick
    }
]

\node[vr] (quest) {
    Meta Quest 3\\
    Unity + OpenXR\\
    Gemelo digital
};

\node[control, right=2cm of quest] (pc) {
    PC de borde\\
    ROS 2 + MoveIt 2\\
    IA + YOLO
};

\node[robot, right=2cm of pc] (ur) {
    Robot UR3\\
    Control físico
};

\node[robot, right=1.5cm of ur] (gripper) {
    Gripper\\
    Producto
};

\node[sensor, below=1.2cm of ur] (camera) {
    Cámara RGB-D\\
    RealSense D405
};

\node[sensor, below=3cm of quest] (loads) {
    Celdas de carga\\
    + HX711
};

\node[sensor, right=1.5cm of loads] (esp) {
    ESP32\\
    Lectura de sensores
};

\node[data, right=1.5cm of esp] (rasp) {
    Raspberry Pi 5\\
    MQTT + PostgreSQL
};

\node[data, right=1.5cm of rasp] (cloud) {
    Nube\\
    Inventario remoto
};

\node[data, right=1.5cm of cloud] (app) {
    Aplicación móvil\\
    Monitoreo
};

\draw[doble] (quest) -- node[above]{WebSocket} (pc);

\draw[doble] (pc) -- node[above]{ROS 2 / RTDE} (ur);

\draw[flecha] (ur) -- (gripper);

\draw[flecha] (camera) -- node[right]{USB 3.0} (pc);

\draw[flecha] (loads) -- (esp);

\draw[flecha] (esp) -- node[above]{MQTT} (rasp);

\draw[doble] (rasp) -- node[above]{HTTPS} (cloud);

\draw[doble] (cloud) -- node[above]{HTTPS} (app);

\draw[doble] (rasp) -- node[right]{MQTT / REST} (pc);

\end{tikzpicture}
}

\caption{Arquitectura general propuesta para el gemelo digital, UR3 e inventario.}

\end{figure}

% ---------------------------------------------------------
\section{Representación tridimensional del almacén}
% ---------------------------------------------------------

Cada celda del almacén tendrá asociada información como:

\begin{itemize}[leftmargin=1.2cm]

    \item Coordenada lógica $(i,j,k)$.
    \item Coordenada física $(x,y,z)$.
    \item Identificador del producto.
    \item Cantidad disponible.
    \item Peso unitario.
    \item Estado de ocupación.
    \item Orientación recomendada del efector final.
    \item Punto de aproximación del robot.

\end{itemize}

Si cada celda tiene dimensiones:

\[
\Delta x,\qquad \Delta y,\qquad \Delta z
\]

y el origen físico del almacén está definido por:

\[
(x_0,y_0,z_0)
\]

la posición central de una celda puede calcularse como:

\[
x=x_0+\left(i+\frac{1}{2}\right)\Delta x
\]

\[
y=y_0+\left(j+\frac{1}{2}\right)\Delta y
\]

\[
z=z_0+\left(k+\frac{1}{2}\right)\Delta z
\]

De esta forma, cuando el usuario selecciona una celda dentro de la
realidad mixta, el sistema puede convertir automáticamente la posición
virtual en una posición física respecto a la base del robot.

Sin embargo, para realizar el agarre correctamente no se utiliza
únicamente $(x,y,z)$, sino una pose completa del efector:

\[
P =
(x,y,z,R_x,R_y,R_z)
\]

donde los últimos tres valores representan su orientación.

% ---------------------------------------------------------
\section{Detección de productos}
% ---------------------------------------------------------

En el efector final se propone colocar una \textbf{cámara RGB-D}.

La imagen RGB será procesada mediante un modelo de visión, por ejemplo
\textbf{YOLO}, para identificar qué producto se encuentra frente al robot.

La información de profundidad permitirá estimar con mayor precisión la
posición del producto antes de realizar el agarre.

El proceso básico será:

\[
\text{Cámara}
\rightarrow
\text{YOLO}
\rightarrow
\text{Identificación}
\rightarrow
\text{Estimación 3D}
\rightarrow
\text{Gripper}
\]

Una cámara RGB convencional podría identificar el objeto, pero no
proporcionaría directamente información de profundidad, por lo que se
selecciona una cámara RGB-D.

% ---------------------------------------------------------
\section{Control del inventario}
% ---------------------------------------------------------

Para mantener actualizado el inventario se propone colocar
\textbf{celdas de carga} en las diferentes posiciones o contenedores del
almacén.

Si una celda contiene un único tipo de producto y el peso individual es:

\[
m_u
\]

mientras que el peso total detectado es:

\[
m_t
\]

se puede estimar la cantidad mediante:

\[
N \approx \frac{m_t}{m_u}
\]

El ESP32 realiza la adquisición de los sensores y envía los valores por
MQTT hacia una Raspberry Pi.

La Raspberry Pi mantiene una base de datos local con información como:

\begin{itemize}[leftmargin=1.2cm]

    \item ID del producto.
    \item Ubicación $(i,j,k)$.
    \item Cantidad.
    \item Peso esperado.
    \item Último movimiento.
    \item Fecha y hora.
    \item Estado de la posición.

\end{itemize}

Cuando el UR3 retira un producto:

\[
N_{\text{nuevo}} = N_{\text{anterior}} - 1
\]

y posteriormente el valor se compara contra la lectura real de la
celda de carga.

Esto permite tener dos fuentes de información:

\begin{enumerate}[leftmargin=1.2cm]

    \item El movimiento registrado por el robot.
    \item El inventario físico obtenido mediante sensores.

\end{enumerate}

% ---------------------------------------------------------
\section{Inteligencia artificial y singularidades}
% ---------------------------------------------------------

La inteligencia artificial no tendrá autoridad directa para mover el
robot.

Primero se utilizará el modelo matemático del UR3 para identificar
configuraciones cercanas a una singularidad.

Una forma de evaluar esta condición es mediante el número de condición
del Jacobiano:

\[
\kappa(J)=
\frac{\sigma_{\max}(J)}
{\sigma_{\min}(J)}
\]

Cuando:

\[
\sigma_{\min}(J)\rightarrow 0
\]

el robot se aproxima a una configuración singular.

El agente de inteligencia artificial recibe información como:

\begin{itemize}[leftmargin=1.2cm]

    \item posición objetivo;
    \item configuración articular;
    \item proximidad a singularidades;
    \item distancia a obstáculos;
    \item límites articulares;
    \item diferentes soluciones de cinemática inversa.

\end{itemize}

Su función será proponer una alternativa, por ejemplo:

\begin{itemize}[leftmargin=1.2cm]

    \item modificar la orientación del efector;
    \item utilizar otra solución de cinemática inversa;
    \item generar un punto intermedio;
    \item aproximarse al producto desde otra dirección.

\end{itemize}

Después, MoveIt vuelve a comprobar matemáticamente la trayectoria.

Por lo tanto:

\[
\text{IA propone}
\rightarrow
\text{MoveIt valida}
\rightarrow
\text{Gemelo ejecuta}
\rightarrow
\text{UR3 ejecuta}
\]

La IA funciona como un sistema de asistencia para encontrar mejores
trayectorias, pero la validación final permanece en algoritmos
determinísticos.

% ---------------------------------------------------------
\section{Diagrama de funcionamiento}
% ---------------------------------------------------------

\begin{figure}[H]
\centering

\begin{tikzpicture}[
    scale=0.72,
    transform shape,
    node distance=0.8cm and 2.8cm,
    iniciofin/.style={
        ellipse,
        draw,
        thick,
        minimum width=3.8cm,
        minimum height=1cm,
        align=center,
        fill=green!15
    },
    proceso/.style={
        rectangle,
        rounded corners,
        draw,
        thick,
        minimum width=5.2cm,
        minimum height=1cm,
        align=center,
        fill=blue!10
    },
    ia/.style={
        rectangle,
        rounded corners,
        draw,
        thick,
        minimum width=5.2cm,
        minimum height=1cm,
        align=center,
        fill=purple!12
    },
    decision/.style={
        diamond,
        draw,
        thick,
        aspect=2.2,
        minimum width=4.2cm,
        minimum height=1.2cm,
        align=center,
        fill=orange!15
    },
    flecha/.style={
        -{Latex[length=2.5mm]},
        thick
    }
]

% Inicio
\node[iniciofin] (inicio) {Inicio};

\node[proceso, below=of inicio] (seleccion) {
    Usuario selecciona producto\\
    o celda $(i,j,k)$ en realidad mixta
};

\node[proceso, below=of seleccion] (objetivo) {
    PC de borde recibe\\
    posición y orientación objetivo
};

\node[proceso, below=of objetivo] (moveit) {
    MoveIt 2 calcula\\
    trayectoria del UR3
};

\node[decision, below=of moveit] (trayectoria) {
    ¿Trayectoria\\
    válida?
};

% Rama de IA
\node[ia, right=3.8cm of trayectoria] (inteligencia) {
    IA analiza singularidad,\\
    colisión o mala configuración
};

\node[proceso, below=of inteligencia] (correccion) {
    IA propone nueva orientación,\\
    pose o punto intermedio
};

% Gemelo digital
\node[proceso, below=of trayectoria] (gemelo) {
    Gemelo digital ejecuta\\
    la trayectoria
};

\node[decision, below=of gemelo] (simulacion) {
    ¿Simulación\\
    correcta?
};

\node[proceso, below=of simulacion] (ejecutar) {
    Usuario presiona\\
    \textbf{Ejecutar}
};

\node[proceso, below=of ejecutar] (validacion) {
    Validación final con estado\\
    actual del UR3 físico
};

\node[proceso, below=of validacion] (robot) {
    UR3 ejecuta\\
    la trayectoria
};

\node[proceso, below=of robot] (vision) {
    Cámara RGB-D + YOLO\\
    verifica el producto
};

\node[proceso, below=of vision] (inventario) {
    Se actualiza el inventario\\
    en la Raspberry Pi
};

\node[proceso, below=of inventario] (nube) {
    Raspberry Pi sincroniza\\
    la información con la nube
};

\node[iniciofin, below=of nube] (fin) {
    Fin / ciclo continuo
};

% Flechas principales
\draw[flecha] (inicio) -- (seleccion);
\draw[flecha] (seleccion) -- (objetivo);
\draw[flecha] (objetivo) -- (moveit);
\draw[flecha] (moveit) -- (trayectoria);

% Trayectoria válida
\draw[flecha]
(trayectoria) --
node[right]{Sí}
(gemelo);

% Trayectoria no válida
\draw[flecha]
(trayectoria) --
node[above]{No}
(inteligencia);

\draw[flecha]
(inteligencia) --
(correccion);

% IA regresa únicamente a MoveIt
\draw[flecha]
(correccion.east)
-- ++(1.2,0)
|- (moveit.east);

% Simulación
\draw[flecha]
(gemelo) --
(simulacion);

\draw[flecha]
(simulacion) --
node[right]{Sí}
(ejecutar);

% Si falla la simulación vuelve a MoveIt
\draw[flecha]
(simulacion.west)
-- ++(-1.7,0)
|- node[pos=0.20,left]{No}
(moveit.west);

% Ejecución física
\draw[flecha] (ejecutar) -- (validacion);
\draw[flecha] (validacion) -- (robot);
\draw[flecha] (robot) -- (vision);
\draw[flecha] (vision) -- (inventario);
\draw[flecha] (inventario) -- (nube);
\draw[flecha] (nube) -- (fin);

\end{tikzpicture}

\caption{Flujo de funcionamiento del sistema de control del UR3 mediante gemelo digital y realidad mixta.}
\label{fig:flujo-sistema}

\end{figure}

% ---------------------------------------------------------
\section{Opciones descartadas y justificación}
% ---------------------------------------------------------

\begin{table}[H]
\centering
\footnotesize
\renewcommand{\arraystretch}{1.4}

\begin{tabularx}{\textwidth}{
    >{\raggedright\arraybackslash}p{3.6cm}
    >{\raggedright\arraybackslash}p{3cm}
    X
}

\toprule
\textbf{Necesidad} &
\textbf{Opción descartada} &
\textbf{Motivo} \\
\midrule

Comunicación VR--robot &
Conectar Quest directamente al UR3 &
Se perdería la capa intermedia de validación de colisiones, cinemática y singularidades. \\

Control del UR3 &
MQTT &
MQTT funciona muy bien para telemetría, pero no es la mejor opción para ejecutar trayectorias del robot. Se selecciona ROS 2 y RTDE. \\

Conexión UR3 &
Wi-Fi &
Puede introducir variaciones de latencia o pérdidas. El robot y la PC se conectarán mediante Ethernet. \\

Procesamiento principal &
Raspberry Pi &
Aunque puede ejecutar tareas básicas, ROS 2, MoveIt, visión e IA simultáneamente requieren mayor capacidad de procesamiento. \\

Visión &
Cámara RGB convencional &
Identifica objetos pero no proporciona directamente información tridimensional. \\

Inventario &
Sensor infrarrojo por celda &
Detecta presencia, pero no determina correctamente la cantidad cuando existen varios productos en la misma posición. \\

Inventario &
Ultrasonido &
La lectura depende de la geometría y posición del producto y puede presentar reflexiones. \\

Inventario &
RFID en todos los productos &
Permitiría identificar individualmente cada pieza, pero aumenta costo y requiere colocar una etiqueta en cada producto. Se considera como mejora futura. \\

Base de datos &
Solamente nube &
Si se pierde Internet, el sistema perdería acceso al estado operativo. Por ello se mantiene una base local y posteriormente se sincroniza. \\

Realidad mixta &
Meta Quest 3S &
Aunque también permite MR, se prefiere Quest 3 por su mejor sistema óptico y sensor de profundidad. \\

Singularidades &
IA como único método &
Una IA puede proponer correcciones, pero no debe reemplazar la evaluación matemática del Jacobiano, límites y colisiones. \\

\bottomrule
\end{tabularx}

\caption{Alternativas descartadas y razones de la toma de decisión.}

\end{table}

% ---------------------------------------------------------
\section{Comunicación con la nube}
% ---------------------------------------------------------

La nube se utilizará exclusivamente para información de alto nivel,
por ejemplo:

\begin{itemize}[leftmargin=1.2cm]

    \item inventario actual;
    \item entradas y salidas;
    \item historial de movimientos;
    \item productos con bajo inventario;
    \item productos más utilizados;
    \item posiciones ocupadas;
    \item alertas;
    \item estadísticas.

\end{itemize}

La nube \textbf{no enviará movimientos directamente al UR3}.

El flujo será:

\[
\text{UR3 / sensores}
\rightarrow
\text{Servidor local}
\rightarrow
\text{Nube}
\rightarrow
\text{Aplicación móvil}
\]

De esta manera, el dueño puede consultar el estado del almacén desde
cualquier lugar sin que la conexión a Internet forme parte del lazo de
control crítico del robot.

% ---------------------------------------------------------
\section{Limitación física del prototipo}
% ---------------------------------------------------------

El UR3 tiene aproximadamente \textbf{500 mm de alcance} y una carga útil
nominal de aproximadamente \textbf{3 kg}.

Por esta razón, la cuadrícula $10\times10\times10$ debe entenderse
principalmente como el modelo lógico del almacén.

El prototipo físico deberá utilizar inicialmente una sección reducida
del almacén ubicada dentro del espacio de trabajo del UR3.

Además, dentro de la carga útil deben considerarse:

\begin{itemize}[leftmargin=1.2cm]

    \item peso del gripper;
    \item peso de la cámara;
    \item soportes;
    \item producto manipulado.

\end{itemize}

Una implementación futura a mayor escala podría utilizar:

\begin{itemize}[leftmargin=1.2cm]

    \item un robot con mayor alcance;
    \item un séptimo eje lineal;
    \item varios robots;
    \item o un robot montado sobre una plataforma móvil.

\end{itemize}

% ---------------------------------------------------------
\section{Regla de seguridad del sistema}
% ---------------------------------------------------------

Una trayectoria correcta dentro del gemelo digital no significa
automáticamente que la misma trayectoria sea segura en el sistema real.

Antes de ejecutar físicamente se debe verificar:

\begin{itemize}[leftmargin=1.2cm]

    \item calibración entre gemelo y robot físico;
    \item posición real del robot;
    \item posición real de los objetos;
    \item TCP del gripper;
    \item carga útil;
    \item límites articulares;
    \item colisiones;
    \item singularidades;
    \item estado de seguridad del UR3.

\end{itemize}

Por lo tanto, se utiliza el principio:

\[
\boxed{
\text{Seleccionar}
\rightarrow
\text{Planear}
\rightarrow
\text{Simular}
\rightarrow
\text{Validar}
\rightarrow
\text{Ejecutar}
}
\]

% ---------------------------------------------------------
\section{Conclusión}
% ---------------------------------------------------------

La propuesta utiliza una arquitectura distribuida en la que cada
dispositivo realiza únicamente la función para la que resulta más
adecuado.

El \textbf{Meta Quest 3} funciona como interfaz de realidad mixta;
la \textbf{PC de borde} realiza la planeación, visión e inteligencia
artificial; el \textbf{UR3} ejecuta únicamente trayectorias previamente
validadas; y la \textbf{Raspberry Pi} administra el inventario local y
la comunicación con la nube.

La información física del almacén se relaciona con una representación
tridimensional parametrizada, permitiendo seleccionar virtualmente un
producto, analizar la trayectoria del robot, corregir posibles
singularidades y posteriormente ejecutar el movimiento sobre el robot
real.

Finalmente, el sistema de inventario mantiene sincronizados el
almacén físico, el gemelo digital, la base de datos y la aplicación
móvil, permitiendo conocer de manera remota qué productos existen,
cuántos quedan y qué movimientos se han realizado.

\end{document}